\begin{textsample}{NEG Dim 3 – global\_north\_2023\_11 – Score -1.00 – global\_north\_2023\_11/91725118.txt}  \label{ex:f3_neg_281}
Rishi Sunak is currently on a two-day trip to the Middle East, where he is visiting the Israeli prime \textbf{minister}, Benjamin Netanyahu, and the country ' s president, Isaac Herzog.
He will also meet others"from across the Middle East", according to Downing Street - including Saudi Arabia."He expressed the deep condolences of the British people for the terrible loss of life that has occurred in Israel.
He underscored the UK ' s firm belief in Israel ' s right to self-defence in accordance with international humanitarian law, as they work to end the threat of Hamas and secure the freedom of hundreds of Israeli hostages."The prime \textbf{minister} thanked Prime \textbf{Minister} Netanyahu for his support for the British nationals who have been taken by Hamas and both leaders agreed to work closely together to secure their freedom."The prime \textbf{minister} welcomed Prime \textbf{Minister} Netanyahu ' s announcement yesterday on opening up aid access to Gaza.
He emphasised the importance of establishing sustained access to get more vital food, water, medicine and fuel into Gaza and to enable British @ @ @ @ @ @ @ @ @ @ underscored the need to prevent any regional escalation in the conflict and the importance of restoring peace and stability to the region."A diplomatic challenge for Sunak Amid the broader messaging around solidarity and sorrow, you can pick out a few more pointed political and diplomatic nuggets from the comments made by the Israeli and British prime \textbf{ministers} at their press conference today.
Benjamin Netanyahu put particular emphasis on the potential length and scale of this conflict with Hamas, saying this will be a"long war"with"ups and downs"that will need"continuous support"from allies.
Rishi Sunak came close to giving public backing to a likely Israeli ground operation into Gaza saying he supports the country ' s right to"go after Hamas"and"to take back hostages".
But he also had something of a veiled warning about the humanitarian situation in Gaza saying"we mourn the loss of every innocent life, civilians of every faith".
With an expectant glance towards his @ @ @ @ @ @ @ @ @ @ Palestinian people are victims of Hamas too."This is a point that has been made by the Israeli side in the last two weeks, but on this occasion, it was not met with any expression of agreement by a neutral-faced Mr Netanyahu.
This isn't the first war Rishi Sunak has dealt with during his short time in office, but the diplomacy here is arguably more complex than with Ukraine.
The challenge for many leaders is finding a balance between showing unwavering support for Israel while recognising the potential humanitarian catastrophe brewing in Gaza and the horrible possibility of conflict spreading across the wider region.
Mr Cleverly posted on social media that he was heading to Egypt, Qatar, and Turkey, where his goals were :"Secure the release of British hostages.
Stop the violence spreading to the region.
Ensure emergency aid can get into Gaza."Poverty-fighting charity Oxfam, which launched an appeal for peace in the region, sided with Mr Sunak on the"safe return of @ @ @ @ @ @ @ @ @ @ call for ceasefire was"wrong".
Katy Chakrabortty, Oxfam ' s head of advocacy said :"The humanitarian situation within Gaza is now reaching catastrophic levels - with more than 3,000 Palestinians dead, what further evidence does he need that any precautions to protect civilians are not working?"The aid allowed represents less than a fifth of that Gaza received every day before the crisis."Wanting Israel to ' win ' is not a helpful statement.
Nobody wins when the scale of human suffering, devastation and pain is so vast and relentless."Hamas officials claimed the hospital blast killed hundreds of people and was caused by an Israeli airstrike - but the Israeli military blamed a misfiring rocket from the Palestinian Islamic Jihad group and released imagery and communications intercepts to support their case.
Spreaker This content is provided by Spreaker, which may be using cookies and other technologies.
To show you this content, we need your permission to use cookies.
You can use the buttons below to @ @ @ @ @ @ @ @ @ @ those cookies just once.
You can change your settings at any time via the Privacy Options.
Unfortunately we have been unable to verify if you have consented to Spreaker cookies.
To view this content you can use the button below to allow Spreaker cookies for this session only.

% matched lemmas: minister
\end{textsample}
